\section{为什么 Transformer 适合生物医学数据}

前一章确定了课程地图，本章从表示层回答共同问题。Transformer 最直接的优势并不是“懂医学”，而是能够在长序列中建立上下文依赖，并用相近的计算骨架接收离散 token、事件、氨基酸或核苷酸。这个统一视角很强，但只有在同时保留各领域语义时才成立。

\subsection{四类序列不是同一种语言}

讲者先请现场学生回答为什么 Transformer 适合 biomedicine，然后把回答落到四个具体对象。clinical note 是医生书写的文本序列；electronic medical record（EMR，电子病历）可以看作患者与医疗系统交互形成的事件序列；protein 是通过肽键连接的氨基酸序列；genome 则是核苷酸碱基序列。它们都存在顺序，却拥有完全不同的“词表”、时间尺度和错误来源。

\lecturefigure{slide-03-why-transformers-biomedicine.jpg}{第一性原理问题：为什么 Transformer 适合生物医学？}{官方录像恢复幻灯片 V003；Stanford CS25 Lecture 17}

\lecturefigure{slide-04-clinical-notes-sequences.jpg}{临床笔记：医生语言形成的文本序列}{官方录像恢复幻灯片 V004；Stanford CS25 Lecture 17}

\begin{knowledgebox}{读图：先区分“序列单位”}
临床笔记中的位置通常表示文本顺序，但事件记录中的位置可能代表时间、就诊或检查；蛋白质中的位置对应 amino acid residue，基因组中的位置对应 nucleotide。共同点是上下文会改变当前位置的意义，差异则是模型必须使用领域适配的 tokenization、位置编码、目标函数和校验方式。
\end{knowledgebox}

电子病历页强调“encounters with the medical system”：患者的诊断、处方、化验和就诊并不只是自然语言段落，而是带有时间、不规则间隔和缺失值的事件流。蛋白质页进一步下沉到分子层级，氨基酸的线性排列会折叠为三维结构，并影响结合位点与功能；因此，序列建模可以提供表示，却不能绕开结构与实验。

\lecturefigure{slide-05-electronic-medical-records.jpg}{电子病历：患者与医疗系统交互的事件序列}{官方录像恢复幻灯片 V005；Stanford CS25 Lecture 17}

\lecturefigure{slide-06-protein-sequences.jpg}{蛋白质：由氨基酸连接形成的生物序列}{官方录像恢复幻灯片 V006；Stanford CS25 Lecture 17}

\teachervoice{讲者把 clinical notes 戏称为“doctor gibberish”，随即改口为 doctor speak。这个玩笑包含真实设计问题：模型也许能生成医学术语，但患者能否理解、医生能否快速核查，是后面 human evaluation 必须单独测量的维度。}

\subsection{基因组与“序列无处不在”}

基因组页展示核苷酸碱基对，随后把临床文本、影像、蛋白质和 DNA 放进同一张总览图。这里的 multimodal（多模态）不是把所有数据强行转成一句话，而是让不同模态保留各自编码，再通过共同表示或跨模态注意力建立联系。比如影像不是天然的一维序列，但可以被切成 patch token；基因表达轨迹也不是语言，却可以作为每个位置的预测目标。

\lecturefigure{slide-07-genome-sequences.jpg}{基因组：核苷酸碱基对构成的长序列}{官方录像恢复幻灯片 V007；Stanford CS25 Lecture 17}

\lecturefigure{slide-08-biomedical-sequences-everywhere.jpg}{生物医学中的序列与模态无处不在}{官方录像恢复幻灯片 V008；Stanford CS25 Lecture 17}

\begin{longtable}{@{}L{0.16\textwidth}L{0.19\textwidth}L{0.22\textwidth}L{0.31\textwidth}@{}}
\toprule
术语 & 序列单位 & 主要上下文 & 不能忽略的领域约束 \\
\midrule
clinical note & subword / word & 病史与当前问题 & 缩写、否定、时间、隐私、临床共识 \\
EMR event & 就诊/诊断/药物/化验 & 患者纵向轨迹 & 不规则时间、缺失、编码体系与选择偏差 \\
protein sequence & amino acid & 局部 motif 与远程残基关系 & 三维折叠、进化保守性、实验功能 \\
genome sequence & nucleotide & 调控元件与长程作用 & 非编码区、细胞类型、染色质环境与测量噪声 \\
multimodal input & text/image/signal token & 模态内与跨模态关系 & 对齐、采样偏差、不同误差尺度 \\
\bottomrule
\end{longtable}

这张术语表还隐含一个重要的数据工程结论：所谓“规模化预训练”在四类数据上并不是同一件事。临床文本的规模受隐私、机构书写习惯和去标识化约束；EMR 事件常因保险、就诊可及性与编码制度产生选择偏差；蛋白质数据库大量序列来自相近物种，随机切分可能让训练集和测试集共享近亲；基因组 track 则集中在少数细胞系与实验协议。若只报告 token 数而不报告来源、去重、家族或个体切分，模型可能通过相似样本记忆获得乐观结果。因此，序列表示之后的第一项工程工作应是定义独立性：患者级、时间级、蛋白质家族级或染色体级切分分别阻断什么泄漏。

\subsection{共同计算骨架与适用边界}

总览页把优势压缩成三个词：sequence modeling、long-range dependencies、scaling。标准 self-attention 对长度为 $n$ 的序列构造两两相似度，因此能直接连接远距离位置；预训练又允许模型利用大量未标注或弱标注数据。但计算和统计优势并不自动带来临床安全、蛋白质功能真实性或变异因果性，后续四个案例正是在补这些缺口。

\lecturefigure{slide-09-transformer-biomedical-fit.jpg}{Transformer 与生物医学数据的匹配点}{官方录像恢复幻灯片 V009；Stanford CS25 Lecture 17}

\begin{importantbox}{统一表示，不等于统一证据}
Transformer 提供的是统一的计算语言：token 表示、上下文交互与可扩展训练。医学结论仍需要领域证据闭环：临床答案由医生与患者评价，蛋白质描述要回到数据库和实验，测序纠错要看 read-level accuracy，调控预测要与真实 genomic tracks 和 variant assays 对照。
\end{importantbox}

\begin{warningbox}{“任何东西都是序列”的误区}
只要把对象排成一列，并不意味着模型学到了正确不变量。EMR 的时间间隔、蛋白质的三维邻近、DNA 的细胞类型依赖都可能被一维顺序隐藏。序列化是建模入口，不是领域结构的替代品。
\end{warningbox}

\subsection{本章小结}

临床文本、医疗事件、蛋白质和基因组共享顺序与上下文依赖，但 token 语义、长度、噪声和验证方式完全不同。Transformer 的价值在于提供可迁移的计算骨架；高质量生物医学系统则必须为每种数据补回领域结构与独立证据。

